% Chapter 3

\chapter{Background and Related Work}
\label{ch:related-work}

This chapter reviews three strands of work that motivate the study.
Section~\ref{sec:background-vqa} considers visual question answering and shortcut
behaviour. Section~\ref{sec:background-rationales} reviews rationale generation and
the distinction between plausibility and faithfulness.
Section~\ref{sec:background-behavioural} examines behavioural evaluation through
interventions and through automatic comparison of generated text.
Section~\ref{sec:background-position} then positions this thesis at their
intersection.

\section{Visual Question Answering}
\label{sec:background-vqa}

\subsection{Benchmarks and shortcut behaviour}
\label{sec:background-shortcuts}

Visual question answering asks a model to answer a question about an image. The
question is written in natural language, and performance is reported as accuracy
over a large evaluation set. That number is necessary, but it can overstate how
much the image contributed, because answers correlate with the words of the
question. \citet{DBLP:conf/cvpr/GoyalKSBP17} give the extreme case: on the first
VQA dataset (VQA v1), answering \enquote{yes} to every question beginning
\enquote{Do you see a\ldots} reaches 87 per cent accuracy without reading the
rest of the question or looking at the image.

Two datasets attacked this problem from different directions. VQA v2 rebalanced
the data question by question. Annotators sought a visually similar second image
on which the same question had a different answer
\citep{DBLP:conf/cvpr/GoyalKSBP17}. Suitable pairs were not always found. For a
successfully contrastive pair, the question alone cannot distinguish the answers.
A model that received only the question fell
from 48.21 to 41.40 accuracy when moved to the balanced evaluation set, which the
authors describe as alleviating the language bias rather than eliminating it.
Visual Question Answering under Changing Priors (VQA-CP) regrouped questions so
that answers common for a question type during training were rare at evaluation
\citep{DBLP:conf/cvpr/AgrawalBPK18}. Every existing model they benchmarked lost
accuracy. The Stacked Attention Network fell from 55.86 on VQA v1 to 26.88 on
VQA-CP v1. The two designs are not
interchangeable. VQA v2 changes what an individual record demands, while VQA-CP
compares models trained on one answer distribution and evaluated on another.

Later work asked whether the remaining shortcuts involve the image as well as the
question. \citet{DBLP:conf/iccv/DancetteCTC21} mined predictive rules from the
training set that combine question words with the labels of detected objects, such
as the presence of a racket, and then collected the validation records on which
every applicable mined rule gives the wrong answer. These counterexamples form
the evaluation subset, and records that no rule matched were set aside. Accuracy
on that subset fell from 63.52
to 33.91 for one common model and from 67.77 to 39.24 for a pretrained
transformer. \citet{DBLP:conf/emnlp/SiMZLL0CWZ22} generalised the idea into nine
shortcut families built from question types, question keywords, annotated key
objects, and their combinations, and reported that accuracy within the training
distribution exceeded accuracy on every one of the nine shifted evaluation sets.
They also found that the methods built to remove such shortcuts did not improve
all nine shifted sets at once.

Together these studies establish that accuracy can coexist with reliance on
regularities the task did not intend, and that those regularities span both the
question and the detected content of the image. Their evidence comes from
rebuilding the data or reselecting the evaluation records. None of them holds a
record fixed and alters its image. \citet{DBLP:conf/cvpr/GoyalKSBP17} do add a
second explanation form, pointing to a similar image on which the answer differs.
Across the work reviewed here, though, no study tests how a generated textual
rationale responds when the visual evidence changes.

\subsection{Visual Commonsense Reasoning}
\label{sec:background-vcr}

VCR was introduced to move visual question answering past recognition towards
inference about people
\citep{DBLP:conf/cvpr/ZellersBFC19}. Its 290,000 questions come from 110,000 movie
scenes and concern activities, social roles, mental states, and likely events
before and after the frame; 38 per cent are explanatory \enquote{why} or
\enquote{how} questions. A formal VCR record supplies the image, a sequence of
object detections carrying a bounding box, a segmentation mask, and a class label,
and a query whose words are either ordinary words or tags referring to one of
those detections. The tags come from a detection model, and the paper notes that
the task is agnostic to how a mask is represented, offering a list of polygons as
one possibility. These tags identify a person or object
only after the reference is resolved against the image.

The dataset poses two subtasks, and both are selections. In \mbox{Q$\rightarrow$A}
the query is the question and the four responses are candidate answers; in
\mbox{QA$\rightarrow$R} the query is the question together with the correct answer
and the four responses are candidate rationales. The authors chose that format
deliberately, \enquote{to avoid the evaluation difficulties of language generation
or captioning tasks where current metrics often prefer incorrect machine-written
text over correct human-written text}
\citep{DBLP:conf/cvpr/ZellersBFC19}. Original VCR therefore never asks a model to
write a rationale. Free-form generation on these records is a later adaptation,
taken up by \citet{DBLP:conf/emnlp/MarasovicBPLSC20} and reviewed in
Section~\ref{sec:background-answer-conditioning}, and it is the adaptation this
thesis uses.

How much of VCR performance rests on the image was tested directly by
\citet{LiGWCNK24}, who removed one input at a time from three vision-language
transformers. Withholding the query cost between 19.7 and 20.5 points on
\mbox{QA$\rightarrow$R}, whereas withholding the image cost between 7.7 and 8.3.
The authors attribute the gap to textual query--response shortcuts, and
illustrate a tag shared between the question and the correct rationale. The
result is bounded to those
architectures and to selection accuracy, and removing the image still cost
something, so it shows uneven rather than absent use of the scene.

VCR thus offers structured references to people and regions, human rationales, and
an established reason to doubt that answer accuracy tracks visual evidence. The
reviewed work measures that doubt by withholding a whole modality from a selection
task. The remaining gap is a joint test of answer behaviour and free-form
rationale behaviour under controlled changes to the same record's image.

\subsection{Descriptions and explicit spatial guidance}
\label{sec:background-descriptions}

A model can also receive scene information through language.
\citet{HakimovS23} test this setup by using captioning and classification models
to convert each image into text; the answering language model receives only that
text. Across five vision and language tasks, they prompted such models with a
small number of examples. The authors describe those results as comparable with
fitted \mbox{vision-language} systems, although the best prompted configuration
trails the fitted comparator on each of the five datasets. Verbal description can
therefore substitute for direct visual input on some tasks.

\citet{DBLP:conf/cvpr/DengCCH25} supply both. They add three kinds of text to
questions over four visual domains: text that matches the image, text corrupted so
that it points to a different answer, and text irrelevant to the question. With
the image retained, corrupted text reduced the accuracy of Molmo-7B-D, an open
\mbox{vision-language} model, from 76.33 to 49.29 on general visual questions. The
proprietary Claude Sonnet moved from 66.88 to 68.17 on the same questions. The
effect of the corrupted text therefore differed across models. Restricting
attention to records where the image alone and the text alone give different
answers, they measure how often the model follows the text. In their certainty
analysis of LLaVA-NeXT-7B and GPT-4o, two further \mbox{vision-language} models,
each model tended to follow whichever channel gave the more certain answer on its
own. Across the full set of models some favoured the image, so this remains a
tendency among the tested systems.

A second way to make the scene explicit is to expose regions rather than describe
them. VCR supplies segmentation regions for its indexed entities
\citep{DBLP:conf/cvpr/ZellersBFC19}, and drawing those regions on the image turns a
tag in the question into a visible pointer. This is related, at the level of
design, to work that makes regions part of the reasoning process:
\citet{DBLP:journals/corr/abs-2403-16999} have a model emit a bounding box for the
region that would help it answer, crop that region, and answer from the original
image and the crop together. The relation is an analogy rather than an
equivalence, since that system predicts its own region while an overlay is drawn
in advance. The ablation of \citet{LiGWCNK24} is a caution against assuming a
marked region will be used. Across the three VCR task variants they report,
replacing every tag in a record with a different tag cost the tested models 0.1 to
1.1 points, while deleting the tags cost 3.9 to 7.9 points. That suggests those
models read the tags as text but made little use of the link to the region.

Two things are established separately. A description can carry enough of a scene
to answer some questions without the image, and it can compete with the image when
the two disagree. Explicit regions can be made part of a reasoning pipeline, though
models tested on VCR have used the link between a tag and its region weakly. What
the reviewed work does not establish is how a fixed description and a simple
polygon overlay relate, together, to both answer behaviour and rationale
sensitivity when the image of a VCR record is altered.

\section{Rationale Generation}
\label{sec:background-rationales}

\subsection{Answer-conditioned rationale models}
\label{sec:background-answer-conditioning}

The answer-conditioned systems reviewed here generate a rationale after receiving
or predicting an answer. The Pointing and Justification model (PJ-X) implements
this design \citep{DBLP:conf/cvpr/ParkHARSDR18}. An answering model predicts an
answer; the predicted answer is embedded and combined with the joint question and
image representation to produce an attention map over the image; the attended visual
feature, the question feature, and the answer embedding are then pooled and passed
to a decoder that writes the justification. The answer is an input to the
explanation rather than a conclusion drawn from it.

Datasets were also developed to train generated rationales and evaluate their
quality. The e-SNLI dataset extends Stanford Natural Language Inference, whose
records pair two sentences and label the relation between
them, with an explanation written by a person \citep{CamburuEtAl18ESNLI}. Their
e-InferSent model was trained in a PredictAndExplain setting: it predicts the label
first, then writes an explanation conditioned on that prediction. The authors graded
the first 100 test examples from one seed by hand, scoring each explanation for
partial as well as full correctness. Among the 80 examples whose label was correct,
the mean correctness score was 34.68 per cent.
\citet{DBLP:conf/emnlp/MarasovicBPLSC20} moved free-form rationalisation to visual
tasks including VCR, and stated their scope plainly: because rationalisation is
hard, they assume the gold answer is given and restrict the study to justifying
that answer. Three crowdworkers then judged whether each rationale supports the
given answer in the context of the image. \citet{KayserEtAl21EVIL} unified this
practice across three datasets, asking annotators \enquote{Given the image and the
question/hypothesis, does the explanation justify the answer?} on a scale with four
points. Their explanation score averages over the examples that were answered
correctly, and their overall score combines it with task accuracy, so an example
with an incorrect answer contributes zero to that overall result.

These studies show that answer-conditioned generators can produce rationales that
people judge as supporting an answer. Surface overlap alone does not capture that
judgement. Their human evaluations measure perceived support or plausibility;
faithfulness requires a separate test. \citet{DBLP:conf/emnlp/MarasovicBPLSC20} also compare each model with a baseline
that receives no image. Human judges then rate how plausible its rationales
appear without visual input, testing whether lexical cues can support a plausible
rationale on their own. When a
generator receives an answer as input, coherence with that answer does not show
whether the rationale reflects the evidence that led to it.

The same dataset therefore supports two distinct tasks. Original VCR asks a
model to pick one rationale from four written by people
\citep{DBLP:conf/cvpr/ZellersBFC19}; free-form generation asks it to compose one.
Selection accuracy and generated text raise different evaluation problems, and the
next subsection separates the properties that such evaluation can and cannot
establish.

\subsection{Plausibility and faithfulness}
\label{sec:background-plausibility}

\citet{DBLP:conf/acl/JacoviG20} separate two properties that explanation research
often runs together. Plausibility is how convincing an interpretation is to a
person. Faithfulness is how accurately it reflects the reasoning process the model
actually used. They observe that the two can come apart, and single out the case
that matters here: where a separate component generates a textual explanation
trained on human explanations, plausibility is the dominating property and there is
no guarantee of faithfulness. \citet{DBLP:journals/corr/abs-2402-04614} restate the
same pair for large language models, defining an explanation as plausible when it
is coherent with human reasoning and understanding, and as faithful when it
accurately represents the reasoning of the underlying model.

The distinction has a direct consequence for the evaluations reviewed above.
\citet{DBLP:conf/acl/JacoviG20} argue that human judgement should not be used to
evaluate faithfulness, because what a person can judge is plausibility, and that
evaluations resting on human gold labels are pushed the same way.
\citet{KayserEtAl21EVIL} accept the point for their own benchmark: they report
human ratings of whether an explanation justifies an answer, and name the
faithfulness of explanations as a question for future work. The human ratings of
Section~\ref{sec:background-answer-conditioning} were therefore designed to
measure plausibility. They do not test faithfulness.

The evaluations above involve four distinct properties. A reader may judge that a rationale
supports an answer. A rationale may resemble a human reference. A model's output
may change when the evidence available to it changes. And there is the process
inside the model that actually produced the answer, which none of the first three
observes directly. The e-SNLI results separate reference overlap from judged
correctness. For the sample described in
Section~\ref{sec:background-answer-conditioning}, the generated explanations came
close to the human annotators' agreement with each other on surface overlap, while
their mean correctness score was 34.68 per cent
\citep{CamburuEtAl18ESNLI}. A fluent and relevant
rationale that readers accept is therefore not yet evidence of faithfulness.

\subsection{Faithfulness of chain-of-thought rationales}
\label{sec:background-cot}

Chain-of-thought prompting supplies a few worked examples in which the answer is
preceded by its intermediate steps, so that the model writes similar steps before
answering \citep{WeiEtAl22CoT}. On the arithmetic tasks, substantial gains appeared
in the models they tested at roughly one hundred billion parameters and above. Two
ablations on GSM8K, a set of grade-school arithmetic word problems, show that the
written sequence matters. Replacing the chain with only the equation to be solved,
or moving the reasoning to after the answer, left performance at approximately the
standard-prompting baseline. Equation-only prompting still helped on simpler one-
or two-step problems. Wei et al. leave open whether emulating the steps
of a human reasoner means that the network is reasoning.

\citet{DBLP:conf/nips/TurpinMPB23} tested whether the written steps report the
reason for the answer. They introduced biasing cues without describing the intended
bias, reordering the options in the examples so that the correct choice was always
the first, or prefixing the question with a stated opinion about the answer. Their
questions came from thirteen tasks in BIG-Bench Hard, a suite of multiple-choice
reasoning problems. Prompted without worked examples, GPT-3.5, one of the two
language models they tested, wrote intermediate steps and scored 59.6 on the
examples whose correct answer contradicted the suggested one. Adding the
suggested answer reduced accuracy from 59.6 to 23.3, a 36.3-point decline. None of
the 234 explanations they inspected from that benchmark mentioned the bias. Of
the unfaithful explanations they annotated, 73 per cent argued for the
answer the cue pointed to, and 15 per cent contained no visible error at all. A generated
explanation can therefore support a cue-driven answer without revealing the cue
that influenced it.

\citet{DBLP:journals/corr/abs-2307-13702} intervened on the written chain itself.
They truncated it before the answer, inserted mistakes into it, paraphrased it, and
replaced it with strings of full stops, then measured whether the final answer
followed. Dependence varied widely across the eight tasks studied, from heavy
reliance on the chain to near indifference to it. Uninformative filler produced no
gain and paraphrasing changed little, which argues against extra computation
alone or the tested wording choices as the explanation in this setup. Their
stated limitation is that they
have no independent access to the model's internal process against which to check
the chain.

This line of work establishes that written reasoning can be steered by cues it never
mentions, and that whether an answer depends on the written chain can be measured by
editing the chain. Both kinds of intervention act on text: on the words of the
prompt, or on the words of the reasoning. Neither shows how a rationale generator
responds when the visual evidence is removed, corrupted, replaced, or transformed.

\section{Behavioural Evaluation of Explanations}
\label{sec:background-behavioural}

\subsection{Interventions as evidence of dependence}
\label{sec:background-interventions}

A common behavioural test removes selected evidence and measures whether the
model's support for its original prediction falls. The ERASER benchmark, short for
Evaluating Rationales And Simple English Reasoning, turns that reasoning
into two measures for text \citep{DeYoungEtAl20ERASER}. Comprehensiveness is the
fall in the probability of the originally predicted class when the selected
rationale tokens are deleted from the input; sufficiency is the difference
between that probability on the full input and on the rationale tokens alone.
The authors are careful about the standing of these measures, calling how best
to measure faithfulness an open question and offering theirs as a first version.
The measures are defined for spans of text selected from the input. Applying the
same logic to text a model writes requires a separate design.

Attention-based explanations prompted a related debate.
\citet{JainWallace19Attention} permuted the attention weights that trained
classifiers place over their input tokens and searched for alternative
distributions that leave the prediction almost unchanged. They found such
distributions readily. They also reported generally weak agreement between learned
attention and importance measured from gradients, for the bidirectional recurrent
encoders they tested, which read each text in both directions.
\citet{DBLP:conf/emnlp/WiegreffeP19} accepted the concern
but questioned the setup. Their adversarial distributions came from a training
procedure of their own, which fits the adversary inside a model while holding its
attention away from an ordinarily trained baseline. They fed those distributions to
a simple model that must use them as fixed guides. On three of the four datasets,
those distributions performed worse as fixed guides than the attention weights
learned by the base classifier; Anemia was the exception. That indicates that the
trained attention carries information about the relation between tokens and
predictions that cannot be freely substituted. Taken together the two papers mark
a boundary: showing that an
alternative exists establishes that the explanation is not unique, which is weaker
than establishing that it is uninformative.

The interventions reviewed in Section~\ref{sec:background-cot} act on generated
reasoning. Across input removal, attention interventions, and chain editing,
changing candidate evidence tests whether an output depends on it. The resulting
evidence remains behavioural: it records the response under the intervention,
while the model's internal route remains unobserved. An unchanged output can also
reflect redundant evidence elsewhere in the input.

\subsection{Automatic comparison of rationale text}
\label{sec:background-metrics}

Automatic evaluation of generated text often compares a candidate with a
reference. BLEU counts how many of a candidate's short word sequences, its
n-grams, appear in the reference,
clipping each count at the number of times it occurs there, combining the
precisions for several lengths as a geometric mean, and penalising candidates that
are too short \citep{DBLP:conf/acl/PapineniRWZ02}. Papineni et al. state that BLEU needs
to agree with human judgement only when averaged over a corpus, and they give the
example of a system that writes \enquote{East Asian
economy} being penalised heavily because the references happen to read
\enquote{economy of East Asia}. The rationales evaluated here are individual
short texts. BLEU was designed for corpus averages, so these scores require
caution.

Two later measures relax the matching. BLEURT scores a candidate against a
reference using representations from BERT, a pretrained language model. It was
trained first on millions of synthetic pairs built by perturbing sentences, and
then fitted to human ratings from machine translation evaluations
\citep{DBLP:conf/acl/SellamDP20}. BERTScore compares contextual token
representations directly, without a learned quality regression. It embeds both
sentences in context, pairs each token with the most similar token in the other
sentence by cosine similarity, and combines the resulting precision and recall
into an F1 score \citep{DBLP:conf/iclr/ZhangKWWA20}. That makes it tolerant of
rewording, and better than surface overlap at the harder case of sentences that
share their words but not their meaning. On pairs built by swapping words within a
sentence, BERTScore's F1
obtained 0.685 for the area under the receiver operating characteristic curve,
which measures how well the score separates paraphrases from non-paraphrases;
BLEU obtained 0.527. This improved on BLEU and remained well below perfect
separation.

In the e-SNLI setting of Section~\ref{sec:background-answer-conditioning}, the
generated explanations scored 22.40 on BLEU, while the third human explanation
scored 22.51 against the other two \citep{CamburuEtAl18ESNLI}. The authors
concluded that the measure was not reliable for the task and turned to manual
grading, which gave a mean correctness score of 34.68 per cent over the 80
correctly labelled examples.
\citet{KayserEtAl21EVIL} compared automatic scores against human judgements of
whether an explanation justifies an answer, over three datasets and four models.
For the examples the models answered correctly, the Spearman rank correlations with
those judgements were 0.216 for BLEU-4, the four-gram variant, 0.248 for BLEURT,
and 0.293 for BERTScore. On the VCR portion alone they were 0.038, 0.128, and
0.138. The source does not mark the VCR BLEU-4 coefficient as meeting its
$p<0.001$ criterion. Similarity to a human
reference is therefore weak evidence about the quality of a rationale on this kind
of data.

Similarity to a reference and change in a model's own output are different
quantities. A rationale may sit far from any human reference and still be the same
text before and after the image changes; it may sit close to a reference in both
conditions and still be different text. Comparing a model's two outputs for one
record avoids relying on a human reference, but it inherits the limits of the
similarity measure, since an embedding comparison still compresses negation, mistaken
references, and changes in causal structure into a single number. Such a number
needs comparison points before it can be read.

% The BERT and ROC explanations lengthened this subsection by four lines, which
% carried its closing line onto the next page above the section heading. Lending
% this page one line keeps the paragraph whole and restores the page break the
% previous pass established.
\enlargethispage{\baselineskip}

\section{Positioning of This Thesis}
\label{sec:background-position}

This thesis addresses the gap identified across the three strands reviewed above
by jointly evaluating Stage~1 answer behaviour and Stage~2 free-form rationale
behaviour on the same VCR records under the same controlled changes to visual
evidence. In the primary Stage~2 analysis, the supplied gold answer remains
fixed while the image condition changes. The four separately adapted
configurations vary the use of descriptions and polygon overlays. This places
the two stages under one intervention framework while keeping their measures
separate.

The three strands motivate different parts of this design. Visual question
answering motivates the controlled image changes and the analysis of answer
behaviour. Rationale generation motivates the Stage~2 setup in which the
supplied answer is fixed. Behavioural evaluation motivates comparing source and
intervention outputs for the same records.
Table~\ref{tab:background-positioning} summarises the three literature strands,
the remaining gap, and the corresponding design response.

\begin{table}[ht]
\centering
\caption[How the reviewed literature motivates this thesis]{How the reviewed
literature motivates this thesis. Each row summarises one strand and the
corresponding design response.}
\label{tab:background-positioning}
\begin{tabular}{L{3.0cm}L{4.8cm}L{5.7cm}}
\toprule
Strand & What prior work established & Remaining gap and response in this thesis \\
\midrule
Visual question answering &
Accuracy can coexist with reliance on regularities in the question, in detected
image content, and in their combination; on VCR, withholding the query costs
selection accuracy more than withholding the image &
Prior work mainly evaluates answer or rationale selection through redesigned
datasets, shifted splits, selected records, or removal of a whole modality. This
thesis keeps the evaluated question and choices fixed, varies the image
condition, and records answer behaviour across the resulting conditions \\
\addlinespace
Rationale generation &
Generators conditioned on an answer produce text readers accept as justifying it;
plausibility and faithfulness are distinct, and written reasoning can be steered by
a cue it never mentions &
Acceptance by a reader does not show that a rationale reflects the evidence behind
the answer. This thesis generates rationales for the same records. In the primary
Stage~2 analysis, the supplied answer remains fixed while the image condition
changes \\
\addlinespace
Behavioural evaluation &
Altering candidate evidence and observing the output tests dependence; similarity
to a human reference is a different quantity and correlates weakly with human
judgement of rationales &
The studies reviewed here intervene on input spans, attention, or generated
reasoning. This thesis applies the same conceptual image conditions at both
stages and compares outputs for the same records \\
\bottomrule
\end{tabular}
\end{table}

